\documentclass[../../../include/open-logic-section]{subfiles}

\begin{document}
\olfileid[hi]{sth}{ord-arithmetic}{intro}

\olsection{परिचय}

\olref[ordinals][]{chap} में हमने क्रमसूचक संख्याओं का सिद्धांत विकसित किया।
\olref[spine][]{chap} में हमने देखा कि क्रमसूचकों को ऐसी रीढ़ माना जा सकता
है जिसके चारों ओर शेष पदानुक्रम निर्मित होता है। किंतु क्रमसूचकों की यही
एकमात्र भूमिका नहीं। क्रमसूचक अंकगणित करने का कार्य भी है।

\olref[ordinals][idea]{sec} में $\omega$, $\omega+1$ और
$\omega+\omega$ की बात करते समय हम इसका संकेत पहले ही दे चुके थे। उस समय
हमने अनौपचारिक रूप से बात की; अब उसे ठीक ढंग से विस्तृत करने का समय है। फिर
भी उल्लेख कर दें कि इस अध्याय में दर्शन बहुत कम है; केवल तकनीकी विकास और यह
(कुछ) रोचक प्रेक्षण कि एक ही काम दो भिन्न ढंग से किया जा सकता है।

\end{document}
